\documentclass[11pt,twocolumn]{article}

\usepackage{epsfig}
\usepackage{graphicx}
\usepackage{graphpap}
\usepackage{amsmath}
\usepackage[latin1]{inputenc}
\usepackage[spanish]{babel}

\topmargin -2.5 cm
\textheight 9.5in
\oddsidemargin -1cm
%\evensidemargin -1cm
\textwidth 18cm
%opening
\title{}
\author{}
\date{}

\newcommand{\abre}{\textquestiondown}

\begin{document}
\pagestyle{empty}
\sffamily
\twocolumn[\begin{minipage}{7 cm}{\begin{center}
\includegraphics[scale=0.8]{escudo}\end{center}}\end{minipage}
\begin{minipage}{10cm}{\vspace*{0.2cm}\Large\textbf{Universidad de Antioquia}\\[2mm] 
Facultad de Ciencias Exactas y Naturales \\[2mm] Instituto de F\'isica\\[2mm] 
F\'isica II, taller 6\\[2mm] Capacitancia}
\end{minipage}

\vspace{0.5cm}
\hrulefill 
\vspace{1cm}]

\footnote{Las figuras han sido tomadas en su gran mayoria de Physics For Scientist and Engineers 6E By Serway and Jewett}

\begin{enumerate}

%ejercicio
\item en redacci�n...



\end{enumerate}
\end{document}
